\documentclass[11pt]{article}
\usepackage{../homework}

\profname{Dr. Michael Martens}
\assignnum{Homework 03}
\classcode{PHYS 313 Thermodynamics}
\duedate{18 September 2026}

\begin{document}
\begin{problem}{1. Three particles on a lattice}\hfill
\begin{enumerate}[label=(\alph*)]
  \item Unique arrangements.
        \begin{answer}
          We have \( 6 \) slots and \( 3 \) particles, so the total arrangements is \(\binom{6}{3}\).
        \end{answer}
  \item Multiplicity function.
        \begin{answer}
        We just enumerate each possible microstate, \( g(0) = 2, g(\varepsilon_0) = 8, g(2\varepsilon_0) = 10, g(3\varepsilon_0) = 0 \).
        \end{answer}
  \item Partition function.
        \begin{answer}
        \( Z(\tau) = \sum\limits_s \exp(-\varepsilon_s/\tau) = 2 + 8\exp(-\varepsilon_0/\tau) + 10\exp(-2\varepsilon_0/\tau)\).
        \end{answer}
  \item Probability of each energy.
        \begin{answer}
          \( \mathbb{P}(U = \varepsilon_s) = \sum\limits_s \frac{\exp(-\varepsilon_{s}/\tau)}{Z(\tau)} \).
          So, we have
          \begin{equation}
          \label{eq:1}
            \begin{aligned}
              &\mathbb{P}(U=0) = \frac{1}{1 + 4\exp(-\varepsilon_{0}/\tau) + 5\exp(-2\varepsilon_0/\tau)}, \\[1.5ex]
              &\mathbb{P}(U= \varepsilon_0) = \frac{4\exp(-\varepsilon_{0}/\tau)}{1 + 4\exp(-\varepsilon_0/\tau) + 5\exp(-2\varepsilon_0/\tau)} \\[1.5ex]
              &\mathbb{P}(U=2\varepsilon_0) = \frac{5\exp(-2\varepsilon_{0}/\tau)}{1 + 4\exp(-\varepsilon_0/\tau) + 5\exp(-2\varepsilon_0/\tau)} \\[1.5ex]
                &\mathbb{P}(U=3\varepsilon_0) = 0.
            \end{aligned}
          \end{equation}
        \end{answer}
        \item Expectation value of energy \( U = \langle\varepsilon\rangle \).
        \begin{answer}
          We calculate \(\langle\varepsilon\rangle = \sum\limits_s \varepsilon_s \mathbb{P}(U=\varepsilon_s)\).
          \begin{equation}
          \label{eq:3}
          \begin{aligned}
            \langle\varepsilon\rangle &= \frac{4\varepsilon_0\exp(-\varepsilon_0/\tau) + 10\varepsilon_0\exp(-2\varepsilon_0/\tau)}{1 + 4\exp(-\varepsilon_0/\tau) + 5\exp(-2\varepsilon_0/\tau)} \\[1.5ex]
            &= \frac{2\varepsilon_{0}\exp(-\varepsilon_{0}/\tau) \left( 2 + 5\exp(-\varepsilon_0/\tau) \right)}{1 + 4\exp(-\varepsilon_0/\tau) + 5\exp(-2\varepsilon_0/\tau)}.
          \end{aligned}
          \end{equation}
        \end{answer}
        \item Heat capacity \( C_V \).
        \begin{answer}
          The heat capacity for a system of constant volume is defined as \( C_V = \frac{\partial \langle\varepsilon\rangle}{\partial \tau}  \).
          Lets define the variable \( \xi = \exp(-\varepsilon_0/\tau) \) for convenience.
          Hence,
          \begin{equation}
          \label{eq:1exen}
            \langle\varepsilon\rangle = \frac{2\varepsilon_0\xi (2 + 5\xi)}{1 + 4\xi + 5\xi^2},
          \end{equation}
          so that
          \begin{equation}
          \label{eq:5}
          \begin{aligned}
            C_V &= \frac{(1 + 4\xi + 5\xi^2)(4\varepsilon_0 + 20\varepsilon_0\xi) - 2\varepsilon_0\xi(2+5\xi)(4+10\xi)}{(1+4\xi + 5\xi^2)^2} \frac{\partial \xi}{\partial \tau} \\[1.5ex]
                &= \frac{(1+4\xi + 5\xi^2)(4\varepsilon_0 + 20\varepsilon_0\xi) - 2\varepsilon_0\xi(8 + 40\xi + 50\xi^2)}{(1+4\xi + 5\xi^2)^2} \frac{\partial \xi}{\partial \tau} \\[1.5ex]
                &= \frac{4\varepsilon_{0} \left( (1 + 4\xi + 5\xi^2)(1 + 5\xi) - (4\xi + 20\xi^2 + 25\xi^3) \right)}{(1+4\xi + 5 \xi^2)^2} \frac{\partial \xi}{\partial \tau} \\[1.5ex]
                &= \frac{4\varepsilon_0(1 + 4\xi + 5\xi^2 + 5\xi + 20\xi^2 + 25\xi^3 - 4\xi - 20\xi^2  -25 \xi^3)}{(1 + 4 \xi + 5\xi^2)^2} \frac{\partial \xi}{\partial \tau} \\[1.5ex]
                &= \frac{4\varepsilon_{0}(1 + 5\xi + 5\xi^2)}{(1+4\xi + 5\xi^2)^2} \frac{\partial \xi}{\partial \tau} \\[1.5ex]
                &= \frac{4\varepsilon_{0}(1 + 5\xi + 5\xi^2)}{(1+4\xi + 5\xi^2)^2} (\frac{\varepsilon_{0}}{\tau^2}\exp(-\varepsilon_0/\tau)) \\[1.5ex]
                &= \frac{4\varepsilon_{0}^2\xi ( 1 + 5\xi + 5\xi^2 )}{\tau^2(1 + 4\xi + 5\xi^2)^2}.
          \end{aligned}
          \end{equation}
        \end{answer}
  \item Show that when \(\xi = \frac{1}{2}\), then \( \langle \varepsilon \rangle = \frac{18}{17}\varepsilon_0 \).
        \begin{proof}
          From \cref{eq:1exen} we have that
          \begin{equation}
          \label{eq:7}
          \langle \varepsilon \rangle\eval_{\xi=1/2} = \frac{\varepsilon_{0}(2+5/2)}{1 + 2 + 5/4} = \frac{9/2}{17/4}\varepsilon_0 = \frac{18}{17}\varepsilon_0.
          \end{equation}
        \end{proof}
  \item Show that when \(\xi = \frac{1}{2}\), then \( C_V = 0.526(\varepsilon_0/\tau)^2\).

        \begin{proof}
          From \cref{eq:5},
          \begin{equation}
          \label{eq:6}
          \begin{aligned}
            C_V \eval_{\xi = 1/2} &= \frac{2\varepsilon_0^2 (1 + 5/2 + 5/4)}{\tau^2(1 + 2 + 5/4)^2} = \frac{(19/2) \varepsilon_{0}^2}{\tau^2(17/4)^2} = 0.526(\varepsilon_0/\tau)^2.
          \end{aligned}
          \end{equation}
          Solving for \(\varepsilon_{0}/\tau\), we have that \( \xi = \exp(-\varepsilon_0/\tau) = 1/2 \), so \( -\log(2) = -\varepsilon_0/\tau \), which gives \( C_V = (\log(2))^2(0.526) = 0.2527 \), as desired.
        \end{proof}
\end{enumerate}
\end{problem}
\begin{problem}{2. Quantum harmonic oscillator}
  Energy levels \(\varepsilon_s = s\hbar \omega\), \( s \in \mathbb{N}_0 \), system in contact with reservoir at temperature \(\tau\).
  Define \( x = \exp(-\hbar\omega/\tau) \) so that \( x^s = \exp(-s\hbar\omega/\tau) \), and for \( x < 1 \),
  \begin{equation}
  \label{eq:8}
  \sum\limits_{s=0}^{\infty} x^s = \frac{1}{1-x},
\end{equation}
and
\begin{equation}
\label{eq:9}
\sum\limits_{s = 0}^{\infty} sx^s = x \frac{\mathrm{d} }{\mathrm{d} x}\left( \sum\limits_{s=0}^{\infty} x^s \right).
\end{equation}
\begin{enumerate}[label=(\alph*)]
  \item Show \( Z(\tau) = \frac{1}{1-\exp(-\hbar\omega/\tau)} \).
        \begin{proof}
          \( Z(\tau) = \sum\limits_s \exp(-\varepsilon_s/\tau) \), so, from \cref{eq:8},
          \begin{equation}
          \label{eq:10}
          \begin{aligned}
            Z(\tau) &= \sum\limits_{s=0}^{\infty} x^{s} = \frac{1}{1-x} = \frac{1}{1- \exp(-\hbar\omega/\tau)}.
          \end{aligned}
          \end{equation}
        \end{proof}
  \item Show \(\langle\varepsilon\rangle = \frac{\hbar\omega}{\exp(\hbar\omega/\tau) - 1}\).
        \begin{proof}
          \( \langle\varepsilon\rangle = \sum\limits_s \varepsilon_s \mathbb{P}(U=\varepsilon_s) \), so
          \begin{equation}
          \label{eq:11}
          \begin{aligned}
            \langle\varepsilon\rangle &= \sum\limits_{s=0}^{\infty} s\hbar\omega \frac{x^{s}}{Z(\tau)} = \sum\limits_{s=0}^{\infty} s\hbar\omega x^s(1-x) \\
                &= \hbar\omega(1-x) \sum\limits_{s=0}^{\infty} sx^s = \frac{\hbar\omega(1-x)x}{(1-x)^2} = \frac{\hbar\omega x}{(1-x)} \\
            &= \frac{\hbar\omega}{(1-x)/x} = \frac{\hbar\omega}{(x^{-1} - 1)} = \frac{\hbar\omega}{\exp(\hbar\omega/\tau) - 1}.
          \end{aligned}
          \end{equation}
        \end{proof}
  \item Find \( C_V \), show that as \(\tau \rightarrow \infty\), \( C \rightarrow 1 \).
        \begin{proof}
          Here,
          \begin{equation}
          \label{eq:12}
          \begin{aligned}
            C_V &= \frac{\partial \langle\varepsilon\rangle}{\partial \tau} = -\hbar\omega(\exp(\hbar\omega/\tau)-1)^{-2}(- \frac{\hbar\omega}{\tau^2} \exp(\hbar\omega/\tau)) \\
            &= \frac{\hbar^2\omega^2 \exp(\hbar\omega/\tau)}{\tau^2(e^{\hbar\omega/\tau} - 1)^2}.
          \end{aligned}
          \end{equation}
        \end{proof}
\end{enumerate}
\end{problem}

\begin{problem}{3}
  System of two spins, magnetic moment \( m \) in thermal contact with reservoir at temperature \(\tau\).
  Spins in magnetic field \( \mathbf{B} \) with amplitude \( B \), can be oriented up or down.
  Energy of up is \( -m B \) and energy down is \( m B \).
  The two spins interact, with energy \( 0 \) is spins are in the same direction, and energy \( J \) if opposite.
  \begin{enumerate}[label=(\alph*)]
    \item Thermal average magnetic moment \(\langle m \rangle\) of the system.
          \begin{answer}
            \(\langle m \rangle = \sum\limits_s m_s \mathbb{P}(U=\varepsilon_s) = \sum\limits_s \frac{m_{s}\exp(-\varepsilon_s/\tau)}{Z(\tau)}\).

            Our partition function is
            \begin{equation}
            \label{eq:13}
              Z(\tau) = 2\exp(-J/\tau) + \exp(2mB/\tau) + \exp(-2mB/\tau),
            \end{equation}
            and thus
            \begin{equation}
            \label{eq:14}
            \begin{aligned}
              \mathbb{P}(U=-J) &= \frac{2\exp(-J/\tau)}{2\exp(-J/\tau) + \exp(2mB/\tau) + \exp(-2mB/\tau)}, \\[1.5ex]
              \mathbb{P}(U=2mB) &= \frac{\exp(-2mB/\tau)}{2\exp(-J/\tau) + \exp(2mB/\tau) + \exp(-2mB/\tau)}, \\[1.5ex]
              \mathbb{P}(U=-2mB) &= \frac{\exp(2mB/\tau)}{2\exp(-J/\tau) + \exp(2mB/\tau) + \exp(-2mB/\tau)}.
            \end{aligned}
            \end{equation}
Now we evaluate the thermal average of the magnetic moment,
          \begin{equation}
          \label{eq:15}
          \begin{aligned}
            \langle m \rangle &= \sum_s m_s \mathbb{P}(s) \\
            &= \frac{(0) \cdot 2\exp(-J/\tau) + (-2m)\exp(-2mB/\tau) + (2m)\exp(2mB/\tau)}{2\exp(-J/\tau) + \exp(2mB/\tau) + \exp(-2mB/\tau)} \\[1.5ex]
            &= \frac{2m\left( \exp(2mB/\tau) - \exp(-2mB/\tau) \right)}{2\exp(-J/\tau) + \exp(2mB/\tau) + \exp(-2mB/\tau)} \\[1.5ex]
            &= \frac{2m\sinh(2mB/\tau)}{\exp(-J/\tau) + \cosh(2mB/\tau)}.
          \end{aligned}
          \end{equation}
        \end{answer}
    \item Find \( \langle m \rangle \) to first order assuming \( mB \ll \tau \). Evaluate \( \langle m \rangle \) in the limits \( J \rightarrow \infty \), \( J \rightarrow -\infty \).
        \begin{answer}
          Expanding \cref{eq:15} in powers of \( mB/\tau \ll 1 \), we use \( \sinh(x) \approx x \) and \( \cosh(x) \approx 1 \),
          \begin{equation}
          \label{eq:16}
          \langle m \rangle \approx \frac{2m (2mB/\tau)}{\exp(-J/\tau) + 1} = \frac{4m^2 B}{\tau \left(1 + \exp(-J/\tau)\right)}.
          \end{equation}
          Evaluating the two asymptotic limits,
          \begin{itemize}
                  \item \( J \rightarrow \infty \)
            \begin{equation}
              \lim_{J \to \infty} \langle m \rangle = \frac{4m^2 B}{\tau}.
            \end{equation}
                  \item \( J \rightarrow -\infty \)
            \begin{equation}
              \lim_{J \to -\infty} \langle m \rangle = 0.
            \end{equation}
          \end{itemize}
        \end{answer}
  \end{enumerate}
\end{problem}

\begin{problem}{4. The zipper model}\hfill
  \begin{enumerate}[label=(\alph*)]
    \item Show that the partition function can be summed in the form
    \begin{equation}
      Z = \frac{1 - \exp(-(N+1)\varepsilon/\tau)}{1 - \exp(-\varepsilon/\tau)}.
    \end{equation}
    \begin{proof}
      A state of the zipper is uniquely specified by the number of open links \( s \in \{0, 1, \dots, N\} \). The energy of a state with \( s \) open links is \( \varepsilon_s = s\varepsilon \). Defining \( y \equiv \exp(-\varepsilon/\tau) \), the partition function is the finite geometric series,
      \begin{equation}
      \label{eq:17}
        Z(\tau) = \sum_{s=0}^N \exp(-s\varepsilon/\tau) = \sum_{s=0}^N y^s = \frac{1 - y^{N+1}}{1 - y} = \frac{1 - \exp(-(N+1)\varepsilon/\tau)}{1 - \exp(-\varepsilon/\tau)}.
      \end{equation}
    \end{proof}
    \item In the limit \( \varepsilon \gg \tau \), find the average number of open links.
    \begin{answer}
      The average number of open links \( \langle s \rangle \) is
      \begin{equation}
      \label{eq:18}
        \langle s \rangle = \frac{1}{Z} \sum_{s=0}^N s y^s = \frac{y}{Z} \frac{\partial Z}{\partial y} = y \frac{\partial \log Z}{\partial y}.
      \end{equation}
      In the low-temperature limit \( \varepsilon \gg \tau \), we have \( y = \exp(-\varepsilon/\tau) \ll 1 \). Neglecting terms of order \( y^2 \) and higher,
      \begin{equation}
      \label{eq:19}
        Z \approx 1 + y = 1 + \exp(-\varepsilon/\tau).
      \end{equation}
      Thus, to first order,
      \begin{equation}
        \langle s \rangle \approx \frac{0 \cdot 1 + 1 \cdot y}{1 + y} \approx y = \exp(-\varepsilon/\tau).
      \end{equation}
    \end{answer}
  \end{enumerate}
\end{problem}

\begin{problem}{5. A hopping ion in graphene}\hfill
  \begin{enumerate}[label=(\alph*)]
    \item Dipole moments \( \mathbf{p}_1, \mathbf{p}_2, \mathbf{p}_3 \).
    \begin{answer}
      With the negative ion \( -e \) at the origin, the electric dipole moment is \( \mathbf{p} = q\mathbf{r} = e\mathbf{r} \). The three equilibrium positions are located at distance \( a \) at angles \( 210^\circ \), \( 330^\circ \), and \( 90^\circ \),
      \begin{equation}
      \label{eq:20}
      \begin{aligned}
        \mathbf{p}_1 &= ea \left(\cos 210^\circ \hat{\mathbf{x}} + \sin 210^\circ \hat{\mathbf{y}}\right) = ea \left( -\frac{\sqrt{3}}{2} \hat{\mathbf{x}} - \frac{1}{2} \hat{\mathbf{y}} \right), \\[1.5ex]
        \mathbf{p}_2 &= ea \left(\cos 330^\circ \hat{\mathbf{x}} + \sin 330^\circ \hat{\mathbf{y}}\right) = ea \left( \frac{\sqrt{3}}{2} \hat{\mathbf{x}} - \frac{1}{2} \hat{\mathbf{y}} \right), \\[1.5ex]
        \mathbf{p}_3 &= ea \left(\cos 90^\circ \hat{\mathbf{x}} + \sin 90^\circ \hat{\mathbf{y}}\right) = ea\,\hat{\mathbf{y}}.
      \end{aligned}
      \end{equation}
    \end{answer}
    \item Thermal average dipole moment \( \langle \mathbf{p} \rangle \) in an electric field \( \mathbf{E} = E\hat{\mathbf{y}} \).
    \begin{answer}
      The electrostatic potential energy of dipole \( \mathbf{p}_i \) in field \( \mathbf{E} \) is \( \varepsilon_i = -\mathbf{p}_i \cdot \mathbf{E} = -p_{i,y} E \),
      \begin{equation}
        \varepsilon_1 = \frac{1}{2}eaE, \quad \varepsilon_2 = \frac{1}{2}eaE, \quad \varepsilon_3 = -eaE.
      \end{equation}
      By symmetry between sites 1 and 2, \( \langle p_x \rangle = 0 \). The partition function is
      \begin{equation}
        Z = 2\exp\left(-\frac{eaE}{2\tau}\right) + \exp\left(\frac{eaE}{\tau}\right).
      \end{equation}
      The expectation value of the dipole moment along \( \hat{\mathbf{y}} \) is
      \begin{equation}
      \label{eq:21}
      \begin{aligned}
        \langle \mathbf{p} \rangle = \langle p_y \rangle \hat{\mathbf{y}} &= \frac{\left(-\frac{1}{2}ea\right) 2\exp\left(-\frac{eaE}{2\tau}\right) + (ea)\exp\left(\frac{eaE}{\tau}\right)}{2\exp\left(-\frac{eaE}{2\tau}\right) + \exp\left(\frac{eaE}{\tau}\right)} \hat{\mathbf{y}} \\[1.5ex]
        &= ea \left[ \frac{\exp\left(\frac{eaE}{\tau}\right) - \exp\left(-\frac{eaE}{2\tau}\right)}{\exp\left(\frac{eaE}{\tau}\right) + 2\exp\left(-\frac{eaE}{2\tau}\right)} \right] \hat{\mathbf{y}} = ea \left[ \frac{1 - \exp\left(-\frac{3eaE}{2\tau}\right)}{1 + 2\exp\left(-\frac{3eaE}{2\tau}\right)} \right] \hat{\mathbf{y}}.
      \end{aligned}
      \end{equation}
    \end{answer}
    \item Limits as \( E \rightarrow +\infty \) and \( E \rightarrow -\infty \).
    \begin{answer}\hfill
      \begin{itemize}
              \item \( E \rightarrow \infty \)
        \begin{equation}
          \lim_{E \to +\infty} \langle \mathbf{p} \rangle = ea\,\hat{\mathbf{y}}.
        \end{equation}
              \item \( E \rightarrow -\infty \)
        \begin{equation}
          \lim_{E \to -\infty} \langle \mathbf{p} \rangle = -\frac{1}{2}ea\,\hat{\mathbf{y}}.
        \end{equation}
      \end{itemize}
    \end{answer}
  \end{enumerate}
\end{problem}

\begin{problem}{6. Calculating free energy}\hfill
  \begin{enumerate}[label=(\alph*)]
    \item Calculate the partition function \( Z \) for this system.
    \begin{answer}
      The energy levels and their multiplicities are \( g_S(0) = 1 \), \( g_S(\varepsilon_0) = 2 \), and \( g_S(2\varepsilon_0) = 1 \). Defining \( \zeta \coloneq \exp(-\varepsilon_0/\tau) \),
      \begin{equation}
      \label{eq:22}
        Z = \sum_{l} g_S(\varepsilon_l)\exp(-\varepsilon_l/\tau) = 1 + 2\exp(-\varepsilon_0/\tau) + \exp(-2\varepsilon_0/\tau) = \left( 1 + \exp(-\varepsilon_0/\tau) \right)^2.
      \end{equation}
    \end{answer}
    \item Calculate the energy \( U = \langle \varepsilon \rangle \) of the system.
    \begin{answer}
      Using \( U = \tau^2 \frac{\partial \log Z}{\partial \tau} \), we have
      \begin{equation}
      \label{eq:23}
      \begin{aligned}
        U &= \frac{0\cdot 1 + \varepsilon_0 \cdot 2\exp(-\varepsilon_0/\tau) + 2\varepsilon_0 \exp(-2\varepsilon_0/\tau)}{(1 + \exp(-\varepsilon_0/\tau))^2} \\[1.5ex]
        &= \frac{2\varepsilon_0\exp(-\varepsilon_0/\tau)\left(1 + \exp(-\varepsilon_0/\tau)\right)}{(1 + \exp(-\varepsilon_0/\tau))^2} = \frac{2\varepsilon_0\exp(-\varepsilon_0/\tau)}{1 + \exp(-\varepsilon_0/\tau)} = \frac{2\varepsilon_0}{\exp(\varepsilon_0/\tau) + 1}.
      \end{aligned}
      \end{equation}
    \end{answer}
    \item Calculate the entropy of the combined system \( S + R \).
    \begin{answer}
      The multiplicity of the reservoir with energy \( U_0 - \varepsilon_l \) can be expanded around \( U_0 \) using \( \sigma_R(U_0 - \varepsilon_l) \approx \sigma_0 - \frac{\varepsilon_l}{\tau} \), so that
      \begin{equation}
        g_R(U_0 - \varepsilon_l) = \exp(\sigma_R(U_0 - \varepsilon_l)) \approx \exp(\sigma_0)\exp(-\varepsilon_l/\tau) = g_R(U_0)\exp(-\varepsilon_l/\tau).
      \end{equation}
      The total multiplicity \( g_{\text{tot}} \) is
      \begin{equation}
      \label{eq:24}
        g_{\text{tot}} = \sum_l g_S(\varepsilon_l) g_R(U_0 - \varepsilon_l) = g_R(U_0) \sum_l g_S(\varepsilon_l)\exp(-\varepsilon_l/\tau) = g_R(U_0) Z.
      \end{equation}
      Taking the logarithm yields the total entropy
      \begin{equation}
      \label{eq:25}
        \sigma_{\text{tot}} = \log_{\text{tot}} = \log g_R(U_0) + \log Z = \sigma_0 + \log Z.
      \end{equation}
      Because the reservoir entropy at equilibrium is \( \sigma_R(U_0 - U) = \sigma_0 - U/\tau \), the entropy of system \( S \) alone is
      \begin{equation}
      \label{eq:26}
        \sigma = \sigma_{\text{tot}} - \sigma_R = \log Z + \frac{U}{\tau}.
      \end{equation}
    \end{answer}
    \item Using the results for \( U \) and \( \sigma \), find the free energy \( F = U - \tau\sigma \).
    \begin{answer}
      Substituting \cref{eq:26} directly into the definition of the Helmholtz free energy gives
      \begin{equation}
      \label{eq:27}
        F = U - \tau\sigma = U - \tau \left( \log Z + \frac{U}{\tau} \right) = -\tau \log Z.
      \end{equation}
      Substituting the partition function from \cref{eq:22},
      \begin{equation}
      \label{eq:28}
        F = -\tau \log\left[\left(1 + \exp(-\varepsilon_0/\tau)\right)^2\right] = -2\tau \log\left(1 + \exp(-\varepsilon_0/\tau)\right).
      \end{equation}
      This verifies \( F = -\tau \log Z \).
    \end{answer}
  \end{enumerate}
\end{problem}
\end{document}
